% STATUS: DRAFT
\chapter{Quasi-local algebras and the thermodynamic limit}\label{ch:quasilocal}

\begin{physmotiv}
Statistical mechanics lives in the thermodynamic limit: phase transitions, equilibrium states with infinite entropy, and the very notion of a ``phase'' exist only for infinitely many particles. In this limit $\ee^{-\beta H}$ and $\Tr$ cease to make sense, as does the Hamiltonian $H$ itself as an operator. The quasi-local algebra is the framework that survives: observables are local, states are limits of finite-volume states, and dynamics acts by automorphisms. This is also the template for AQFT, where ``finite volume'' becomes ``bounded spacetime region''.
\end{physmotiv}

\section{The framework}

Let $\Gamma$ be a lattice (say $\Z^d$) with a finite-dimensional algebra $\A_x=M_k(\C)$ (or CAR) at each site. For finite $\Lambda\subset\Gamma$ let $\A_\Lambda=\bigotimes_{x\in\Lambda}\A_x$. Then
\begin{itemize}
\item \textbf{isotony:} $\Lambda_1\subset\Lambda_2\Rightarrow\A_{\Lambda_1}\subset\A_{\Lambda_2}$;
\item \textbf{locality:} $[\A_{\Lambda_1},\A_{\Lambda_2}]=0$ for disjoint $\Lambda_1,\Lambda_2$ (graded for fermions);
\item the \emph{local algebra} $\A_{\rm loc}=\bigcup_\Lambda\A_\Lambda$ and the \emph{quasi-local algebra} $\A=\overline{\A_{\rm loc}}^{\norm\cdot}$.
\end{itemize}
This is a \cstar-algebra, simple, with a unique tracial state, and it is the same object regardless of the Hamiltonian. The Hamiltonian enters as an interaction $\Phi$ assigning to each finite $X\subset\Gamma$ a self-adjoint $\Phi(X)\in\A_X$, with $H_\Lambda=\sum_{X\subset\Lambda}\Phi(X)$.

\section{Dynamics without a Hamiltonian}

In infinite volume the formal sum $H=\sum_X\Phi(X)$ does not define an element of $\A$ (nor an operator in a general representation). Instead, define
\begin{equation}
\alpha_t(a)=\lim_{\Lambda\uparrow\Gamma}\ee^{\ii tH_\Lambda}a\,\ee^{-\ii tH_\Lambda}\qquad(a\in\A_{\rm loc}).
\end{equation}
For interactions that decay sufficiently fast (e.g.\ finite range), this limit exists in norm: this is a consequence of the \emph{Lieb--Robinson bound}, which states that the commutator $\norm{[\alpha_t(a),b]}$ for $a,b$ localized at distance $r$ is exponentially small outside the ``light cone'' $r\lesssim vt$. The limit defines a strongly continuous one-parameter group of automorphisms of $\A$ whose generator, defined on local elements, is the derivation $\delta(a)=\ii\lim[H_\Lambda,a]$.

Conclusion: \emph{dynamics is a one-parameter automorphism group $\alpha_t$, generally not inner}. There is no $H\in\A$ with $\alpha_t(a)=\ee^{\ii tH}a\ee^{-\ii tH}$. Only in a given representation $\pi_\omega$ with an $\alpha$-invariant state does $H_\omega$ exist as a self-adjoint operator (via Stone, \cref{ch:spectral}), with $H_\omega\Omega=0$ for a ground state or equilibrium reference. This is the algebraic form of the QFT lesson that the Hamiltonian depends on the vacuum representation.

\begin{whatwrong}
Different invariant states $\omega$ produce different operators $H_\omega$ on different Hilbert spaces, and the spectrum can differ (ground states have $H_\omega\ge0$; the equilibrium states at $T>0$ have $H_\omega$ unbounded in \emph{both} directions, the same pattern as the modular Hamiltonian). Never speak of ``the'' Hamiltonian of an infinite system without a representation.
\end{whatwrong}

\section{Thermodynamic limit of states}

Given a finite-volume Gibbs state $\omega_\Lambda=\Tr(\ee^{-\beta H_\Lambda}\,\cdot)/Z_\Lambda$, regard $\omega_\Lambda$ as a state on $\A$ (by composing with the conditional expectation onto $\A_\Lambda$, or by inserting trivial boundary conditions). The state space of $\A$ is weak-$^*$ compact (\cref{ch:states}), so a subsequence converges; its limit points are candidate infinite-volume equilibrium states. The following points matter:
\begin{enumerate}
\item The limit points are $\alpha_t$-invariant (stationary) states satisfying the \emph{KMS condition} at inverse temperature $\beta$ (\cref{ch:kms}); this is how one defines equilibrium states directly in infinite volume, with no reference to $\ee^{-\beta H}$.
\item The set of KMS states at fixed $\beta$ is a simplex: each is a unique mixture of \emph{extremal} ones, and the extremal KMS states are the pure phases. Several extremal KMS states (e.g.\ $+$ and $-$ magnetized) at a given $\beta$ signal a first-order phase transition or symmetry breaking; a unique one means a unique phase.
\item Extremal KMS states are \emph{factor states}: the GNS representation $\pi_\omega(\A)''$ is a factor (trivial centre, \cref{ch:factors}). Mixtures of phases have a non-trivial centre, the macroscopic order parameters (\cref{ch:inequiv}). The \emph{type} of this factor is a finer invariant: in the examples we will compute (\cref{ch:connes_class}), it is II$_1$ at $\beta=0$ and type III at generic finite temperature.
\end{enumerate}

\section{Why this matters for later chapters}

The infinite spin chain at infinite or finite temperature is the first example of a type II$_1$ or type III algebra arising from physics: $\pi_\tau(\A)''$ is the hyperfinite II$_1$ factor, while GNS representations of product states $\omega_p$ ($p\neq1/2$) give the Araki--Woods/Powers factors of type III$_\lambda$ (\cref{ch:connes_class}). The same structure appears for a local region of a QFT: replace $\Lambda$ by a spacetime region, the lattice by a net of algebras, and the $\Z^d$ translations by Poincar\'e transformations (\cref{ch:haagkastler}).

\begin{keyidea}
The quasi-local algebra encodes kinematics (locality), dynamics is an automorphism group (not inner), equilibrium states are limits of Gibbs states characterized by the KMS condition, and phases correspond to extremal states and to the centre of the GNS von Neumann algebra.
\end{keyidea}

\section*{Exercises}
\begin{exercise}
For the non-interacting chain $H_\Lambda=\sum_{j\in\Lambda}\frac h2\sigma^z_j$, show $\alpha_t(\sigma^\pm_j)=\ee^{\pm\ii ht}\sigma^\pm_j$, so the limit exists trivially. Compute the Gibbs state limit $\omega_\beta$ and show it is a product state with $\omega_\beta(\sigma^z_j)=-\tanh(\beta h/2)$ (with $\sigma^z=\pm1$ and energy $\pm h/2$). Verify $\omega_\beta\circ\alpha_t=\omega_\beta$.
\end{exercise}
\begin{exercise}
Show the Hamiltonian $H=\frac h2\sum_{j\in\Z}\sigma_j^z$ is not an element of $\A$ and not the norm limit of local Hamiltonians $H_N$ (compute $\norm{H_N-H_M}$). Show that in the representation $\pi_\beta$ the generator of $\alpha_t$ exists as a self-adjoint operator but is not bounded below (consider states with local flips).
\end{exercise}
\begin{exercise}
Using the Ising chain with $H_\Lambda=-J\sum_j\sigma_j^z\sigma_{j+1}^z$ at $T=0$, find the two ground states (all up, all down) as product states, show they are disjoint representations, and show the infinitesimal generator $\delta$ is the same in both.
\end{exercise}
\begin{exercise}
State precisely why the convergence $\omega_\Lambda\to\omega$ need only be weak-$^*$ and not in norm: give an example where $\norm{\omega_\Lambda-\omega_{\Lambda'}}$ stays close to $2$ for all $\Lambda\neq\Lambda'$, e.g.\ product states with different $\theta$ restricted to the full chain.
\end{exercise}
